%% ============================================================
%%  Starlink MHR Routing Simulator — Final Year Project Report
%% ============================================================
\documentclass[12pt,a4paper]{report}

\usepackage[utf8]{inputenc}
\usepackage[T1]{fontenc}
\usepackage{lmodern}
\usepackage[a4paper,top=2.8cm,bottom=2.6cm,left=2.8cm,right=2.5cm,headheight=14pt]{geometry}
\usepackage{amsmath,amssymb}
\usepackage{booktabs}
\usepackage{array}
\usepackage{longtable}
\usepackage{tabularx}
\usepackage{float}
\usepackage{graphicx}
\usepackage[dvipsnames]{xcolor}
\usepackage{fancyhdr}
\usepackage{titlesec}
\usepackage{setspace}
\usepackage{parskip}
\usepackage{microtype}
\usepackage{enumitem}
\usepackage{caption}
\usepackage{subcaption}
\usepackage{cite}
\usepackage{url}
\usepackage{listings}
\usepackage[ruled,linesnumbered,vlined]{algorithm2e}
\usepackage[colorlinks=true,linkcolor=blue!60!black,citecolor=green!50!black,urlcolor=blue!60!black]{hyperref}

\pagestyle{fancy}
\fancyhf{}
\fancyhead[L]{\small\leftmark}
\fancyhead[R]{\small\thepage}
\renewcommand{\headrulewidth}{0.4pt}
\fancypagestyle{plain}{\fancyhf{}\fancyhead[R]{\small\thepage}\renewcommand{\headrulewidth}{0pt}}

\titleformat{\chapter}[display]
  {\normalfont\LARGE\bfseries}{\chaptertitlename\ \thechapter}{8pt}{\LARGE}
\titlespacing*{\chapter}{0pt}{16pt}{12pt}

\setlength{\parindent}{0pt}
\setlength{\parskip}{5pt}
\onehalfspacing

\newcommand{\Pint}{P_{\mathrm{int}}}
\newcommand{\Tprop}{\tau_{\mathrm{prop}}}
\newcommand{\Tser}{T_{\mathrm{s}}}
\newcommand{\Tq}{W_{\mathrm{q}}}
\newcommand{\TE}{\tau_{\mathrm{E2E}}}
\newcommand{\RE}{R_E}
\newcommand{\dth}{d_{\mathrm{th}}}
\newcommand{\thetar}{\theta_r}
\newcommand{\thetas}{\theta_s}
\newcommand{\norm}[1]{\left\|#1\right\|}
\newcommand{\abs}[1]{\left|#1\right|}

%% ============================================================
\begin{document}

%% ── Title page ──────────────────────────────────────────────
\begin{titlepage}
  \centering
  \vspace*{2cm}
  {\Large\bfseries Final Year Project Report\par}
  \vspace{0.5cm}
  \rule{\linewidth}{1pt}
  \vspace{0.4cm}
  {\LARGE\bfseries
A Starlink-Like LEO Multi-Hop Routing Simulator and Empirical Evaluation of the BPP Analytical Model
  \par}
  \vspace{0.4cm}
  \rule{\linewidth}{1pt}
  \vspace{2cm}
  {\large\bfseries Ali Wazni \quad \& \quad Jaafar Mukahal\par}
  \vspace{0.5cm}
  {\large May 2026\par}
  \vfill
  {\small
    Reference paper: Wang, Y., Kishk, M.\,A., \& Alouini, M.-S.\ (2023).\\
    \textit{Reliability Analysis of Multi-hop Routing in Multi-tier LEO Satellite Networks}.\\
    arXiv:2303.02286
  }
\end{titlepage}

%% ── Abstract ────────────────────────────────────────────────
\chapter*{Abstract}
\addcontentsline{toc}{chapter}{Abstract}

This project builds a Starlink-like multi-hop LEO routing simulator from real
Two-Line Element (TLE) data. Starting from publicly available Starlink TLEs, the
simulator generates ten frozen-time network snapshots covering 81 minutes of orbital
movement, each containing 4{,}080 satellites and 250 ground gateways linked by roughly
37{,}000 inter-satellite and gateway-access edges.

Two routing algorithms are implemented and compared: Dijkstra shortest-path routing
and the greedy Wang-style BPP algorithm from Wang, Kishk, and Alouini (2023)
\cite{wang2023reliability}. Both are run on the same real orbital geometry, and their
empirical interruption probabilities are compared against the paper's closed-form BPP
analytical prediction.

At $\dth = 4{,}000$~km on the full constellation, BPP predicts 1.44\% interruption;
the empirical greedy simulator measures 78.57\%---an absolute error of 77 percentage
points. Dijkstra achieves zero interruptions on every tested pair, confirming that
the Walker-delta topology is globally routable and that greedy routing failures are
algorithmic, not structural.

A full-featured 3D browser visualiser built on Three.js accompanies the simulator,
showing real satellite positions, ISL links, animated packets, and a live traffic
simulation with congestion heatmaps, M/D/1 queuing, and a performance dashboard.

\textbf{Keywords:} LEO, Starlink, multi-hop routing, BPP, Walker-delta, TLE,
Dijkstra, interruption probability, M/D/1 queueing, 3D visualisation.

\tableofcontents

%% ============================================================
\chapter{Introduction}
\label{ch:intro}
%% ============================================================

LEO satellite constellations such as Starlink, Kuiper, and OneWeb are deploying tens
of thousands of satellites at 340--600~km altitude to deliver sub-50~ms broadband
globally. Unlike geostationary satellites ($\sim$35{,}786~km, $\sim$600~ms RTT), LEO
satellites can achieve competitive latency for interactive applications.

Because real Starlink routing tables and traffic data are proprietary, simulation from
public TLE data is the main academic tool for studying routing behaviour. Analytical
models such as the BPP model of Wang et al.\ offer closed-form predictions without
simulation, but they assume satellites are placed randomly---a strong simplification
relative to real Walker-delta orbital geometry.

This project asks: \textit{does the BPP analytical prediction hold when tested against
real Starlink TLE-derived geometry?} To answer this, we build a complete simulator that
runs both the paper's own greedy routing algorithm and a Dijkstra baseline on the same
real topology, then compare the results against the closed-form BPP prediction.

%% ============================================================
\chapter{Problem Statement and Objectives}
\label{ch:problem}
%% ============================================================

\section{Problem Statement}

Three sub-problems motivate this project:

\begin{enumerate}
  \item \textbf{Data gap.} Real Starlink routing behaviour is proprietary; TLE data
        provides positions but no routing information.
  \item \textbf{Model validity.} The BPP formula assumes randomly placed satellites.
        Real Starlink follows Walker-delta orbital geometry, which is not random.
  \item \textbf{Algorithm suitability.} The paper's greedy routing algorithm is
        designed for BPP geometry. Whether it succeeds on a real Starlink-like
        topology has not been empirically tested.
\end{enumerate}

\section{Objectives}

\begin{itemize}
  \item Build a Starlink-like LEO routing simulator from real TLE data.
  \item Generate multi-epoch network snapshots (satellites + gateways + ISL/access edges).
  \item Implement the Wang/BPP greedy routing algorithm exactly as specified in
        \cite{wang2023reliability}.
  \item Implement Dijkstra shortest-delay routing as a global-knowledge baseline.
  \item Compare both empirical results against the BPP closed-form analytical prediction.
  \item Measure routing success rate, hop count, end-to-end latency, and delay distribution.
  \item Build a real-time 3D visualiser with live traffic simulation and congestion metrics.
\end{itemize}

%% ============================================================
\chapter{Background}
\label{ch:background}
%% ============================================================

\section{Walker-Delta Constellations}

A Walker-delta constellation is parameterised by the triplet $N/P/F$ (satellites /
planes / phasing factor). Starlink's operational shells use inclinations of
${\sim}43\text{\textdegree}$ and ${\sim}53\text{\textdegree}$, arranged in evenly
spaced orbital planes with $\pm 2$ intra-plane and $\pm 2$ inter-plane
inter-satellite links (ISLs). The regular lane structure provides coverage efficiency
but creates angular anisotropy that affects geographic routing.

\section{Routing Approaches}

\textbf{Greedy geographic routing} selects the next relay based on local geometric
progress toward the destination. It requires no global topology knowledge, but can fail
in \emph{routing voids} where no local relay satisfies the forwarding constraints.

\textbf{Dijkstra shortest-path routing} computes the globally optimal path with full
topology knowledge. It cannot get stuck in voids; the trade-off is the need for a
distributed link-state protocol (e.g., OSPF-TE) to maintain topology information.

\section{The BPP Analytical Model}

Wang et al.\ \cite{wang2023reliability} model each orbital tier as a Binomial Point
Process---$N$ satellites placed independently and uniformly at random on a sphere. They
derive the probability that a valid relay exists at each hop given the routing
constraints (direction, dome angle, distance, and line-of-sight), then chain this across
multiple hops via a Markov-chain formulation to obtain the end-to-end interruption
probability. The model is mathematically tractable and has been validated against
Monte Carlo BPP simulation, but not against real Walker-delta TLE data.

%% ============================================================
\chapter{Mathematical Model}
\label{ch:math}
%% ============================================================

\section{Geometry}

\textbf{Distance} between nodes $u$ and $v$ with ECEF positions $\mathbf{p}_u,\mathbf{p}_v$:
\begin{equation}
  d(u,v) = \norm{\mathbf{p}_u - \mathbf{p}_v}_2
\end{equation}

\textbf{Propagation delay} over link $(u,v)$:
\begin{equation}
  \Tprop(u,v) = \frac{d(u,v)}{c}, \quad c = 299{,}792.458~\text{km/s}
\end{equation}

\textbf{Dome angle} (geocentric angle between two nodes):
\begin{equation}
  \theta_{\mathrm{dome}}(u,v) = \arccos\!\bigl(\hat{\mathbf{p}}_u \cdot \hat{\mathbf{p}}_v\bigr),
  \quad \hat{\mathbf{p}} = \mathbf{p}/\norm{\mathbf{p}}
\end{equation}

\textbf{Direction angle} from current node $c$ toward candidate $v$, relative to destination $d$:
\begin{equation}
  \theta_{\mathrm{dir}}(c,v,d) =
    \arccos\!\left(
      \frac{\mathbf{p}_v-\mathbf{p}_c}{\norm{\mathbf{p}_v-\mathbf{p}_c}}
      \cdot
      \frac{\mathbf{p}_d-\mathbf{p}_c}{\norm{\mathbf{p}_d-\mathbf{p}_c}}
    \right)
\end{equation}

\textbf{Line-of-sight} condition (segment $uv$ does not intersect Earth of radius $\RE = 6{,}378.137$~km):
\begin{equation}
  \min_{t\in[0,1]}\norm{\mathbf{p}_u + t(\mathbf{p}_v-\mathbf{p}_u)}^2 \geq \RE^2
\end{equation}

\section{BPP Routing Constraints}

The Wang/BPP algorithm applies three constraints to filter candidate relays $v$ at each
hop from current node $c$ toward destination $d$:

\begin{align}
  c_1:\quad & \theta_{\mathrm{dir}}(c,v,d) \leq \thetar \label{eq:c1}\\
  c_2:\quad & \theta_{\mathrm{dome}}(c,v)  \geq \thetas  \label{eq:c2}\\
  c_3:\quad & d(c,v) \leq \dth \;\text{and}\; \mathrm{LoS}(c,v) \label{eq:c3}
\end{align}

$c_1$ (direction, \cite{wang2023reliability}~Sec.~III-A) restricts the search to a
forward cone; $c_2$ (dome angle) ensures geometric progress; $c_3$ (distance + LoS)
limits hop length and requires clear sky. The relay minimising
$\theta_{\mathrm{dome}}(v,d)$ among valid candidates is chosen:
\begin{equation}
  v^* = \arg\min_{v\in\mathcal{C}}\,\theta_{\mathrm{dome}}(v,d)
\end{equation}
If $\mathcal{C}=\emptyset$, the packet is \textbf{interrupted}.

\section{BPP Analytical Interruption Probability}

For tier $k$ with $N_k$ randomly placed satellites, the probability that no valid relay
exists in that tier is:
\begin{equation}
  q_k = (1-p_k)^{N_k}
\end{equation}
where $p_k = \Omega_{\mathrm{valid}}^{(k)}/(4\pi R_k^2)$ is the fraction of the sphere
surface satisfying all three constraints. The single-hop interruption probability over
$K$ tiers is:
\begin{equation}
  \Pint^{(1)} = \prod_{k=1}^{K}(1-p_k)^{N_k}
\end{equation}
Multi-hop interruption is derived via a Markov chain over dome-angle states
\cite{wang2023reliability}:
\begin{equation}
  \Pint^{(\mathrm{multi})} = 1 - \bigl[\mathbf{T}^{\bar H}\bigr]_{\mathrm{src}\to\mathrm{dst}}
\end{equation}
where $\mathbf{T}$ is the hop transition matrix and $\bar H$ is the expected hop count.

\section{Dijkstra Path Cost}

\begin{equation}
  C(\pi^*) = \min_{\pi\in\Pi(s,d)}\sum_{i=0}^{H-1}\Tprop(v_i,v_{i+1})
\end{equation}

\section{Delay Model}

\textbf{Serialisation delay} for packet of $L$ bytes on link of capacity $R$ bps:
\begin{equation}
  \Tser = \frac{8L}{R}
\end{equation}

\textbf{M/D/1 queuing delay} (Pollaczek--Khinchine, deterministic service):
\begin{equation}
  \Tq = \frac{\rho}{2(1-\rho)}\,\Tser, \quad \rho = \frac{\text{offered load}}{R}
\end{equation}

\textbf{Total end-to-end delay}:
\begin{equation}
  \TE = \sum_{i=1}^{H}\bigl(\Tprop^{(i)} + \Tser + \Tq^{(i)}\bigr)
  \label{eq:e2e}
\end{equation}

%% ============================================================
\chapter{System Overview and Implementation}
\label{ch:impl}
%% ============================================================

\section{Pipeline}

\begin{center}
\small
\begin{tabular}{c}
\fbox{\parbox{0.80\textwidth}{\centering
\textbf{TLE dataset + ground station CSV}\\[2pt]$\downarrow$\\
\textbf{\texttt{tle\_to\_snapshot.py}} — propagate positions, build ISL/access graph, export CSVs\\[2pt]$\downarrow$\\
\textbf{\texttt{analyze\_mhr\_reliability.py}} — BPP formula + greedy + Dijkstra routing\\[2pt]$\downarrow$\\
\textbf{\texttt{make\_plots.py}} — comparison plots\\[2pt]$\downarrow$\\
\textbf{Browser 3D visualiser} — Three.js, live traffic, performance dashboard
}}
\end{tabular}
\end{center}

\section{Dataset}

The TLE file (\texttt{starlink.tle}) contains 9{,}887 raw records. After filtering for
operational satellites (altitude 400--600~km, eccentricity $<0.01$), 8{,}691 valid TLEs
remain. Shell selection retains the two largest well-separated orbital groups: 2{,}670
satellites at ${\sim}43\text{\textdegree}$ inclination and 1{,}410 at
${\sim}53\text{\textdegree}$, totalling 4{,}080 satellites. The ground station CSV
contains 250 real Starlink gateway locations.

Ten frozen-time snapshots are generated at 9-minute intervals starting at 21:30 UTC,
21 March 2026, using the Skyfield SGP4 propagator.

\begin{table}[H]
\centering
\caption{Mean topology statistics across 10 epochs.}
\label{tab:topology}
\begin{tabular}{lc}
\toprule
\textbf{Quantity} & \textbf{Value} \\
\midrule
Satellites (2 shells)            & 4{,}080 \\
Gateways                         & 250 \\
Total nodes                      & 4{,}330 \\
ISL edges per epoch              & $\approx 7{,}850$ \\
Access edges per epoch           & $\approx 29{,}600$ \\
Total edges per epoch            & $\approx 37{,}450$ \\
Mean ISL length                  & $\approx 1{,}246$~km \\
Mean access link length          & $\approx 1{,}367$~km \\
Largest connected component      & $\geq 99.6\%$ of nodes \\
\bottomrule
\end{tabular}
\end{table}

\textbf{ISL policy:} intra-plane $\pm 2$ neighbours, inter-plane $\pm 2$ adjacent
planes, max length 5{,}000~km; ISLs above $|\mathrm{lat}|>70\text{\textdegree}$ are
disabled (polar seam avoidance).

\section{Routing Algorithms}

\subsection{Dijkstra}
Standard binary min-heap Dijkstra on the ISL/access graph with propagation delay as
edge weight. Gateways are endpoints only---they cannot serve as relay nodes mid-path.

\subsection{BPP Greedy (Algorithm~\ref{alg:bpp})}
At each hop, all $N$ nodes are scanned. Constraints $c_1$, $c_2$, $c_3$
(Eqs.~\ref{eq:c1}--\ref{eq:c3}) are applied. The candidate with smallest dome angle
to destination is chosen. If none exists: interrupted. This matches the paper's
specification exactly, including relaxing angular constraints for gateway--satellite
hops.

\begin{algorithm}[H]
\DontPrintSemicolon
\caption{BPP Greedy Multi-Hop Routing}
\label{alg:bpp}
\KwIn{Source $s$, destination $d$, $(\thetar, \thetas, \dth)$}
\KwOut{Path $\pi$, outcome}
$c \leftarrow s$;\; $\pi \leftarrow [s]$;\; visited $\leftarrow \{s\}$\;
\For{hop $= 1$ \KwTo MaxHops}{
  \lIf{$d(c,d)\leq\dth$ and $\mathrm{LoS}(c,d)$}{append $d$; \Return success}
  $\mathcal{C} \leftarrow \{v\notin\text{visited}:c_1\wedge c_2\wedge c_3\}$\;
  \lIf{$\mathcal{C}=\emptyset$}{\Return interrupted}
  $v^*\leftarrow\arg\min_{v\in\mathcal{C}}\theta_\mathrm{dome}(v,d)$;\;
  $\pi$.append$(v^*)$;\; visited.add$(v^*)$;\; $c\leftarrow v^*$\;
}
\Return interrupted\;
\end{algorithm}

\section{3D Visualiser}

The browser visualiser uses Three.js (WebGL) with:
\begin{itemize}[noitemsep]
  \item \textbf{Instanced rendering:} 4{,}080 satellites in a single
        \texttt{InstancedMesh} draw call; $\sim$37{,}000 edges in a single
        \texttt{LineSegments} draw call with per-vertex colours.
  \item \textbf{Animated packet:} a glowing sphere traversing the route hop by hop.
  \item \textbf{Background traffic:} four intensity presets (Low, Medium, High,
        Stress) generating continuous random gateway-to-gateway flows.
  \item \textbf{Live dashboard:} EMA utilisation $\rho_e$, M/D/1 queuing delay,
        packet loss, goodput, jitter, and a per-link congestion heatmap.
  \item \textbf{Dijkstra vs BPP comparison panel:} sends both routing algorithms
        sequentially and shows results side by side.
\end{itemize}

Gateway access links use a lower capacity (25~Mbps default) than ISL trunks (100~Mbps),
making gateways visible bottlenecks in the congestion heatmap.

\begin{table}[H]
\centering
\caption{Background traffic preset parameters.}
\label{tab:presets}
\begin{tabular}{lcccc}
\toprule
\textbf{Intensity} & \textbf{Rate (pkt/s)} & \textbf{Demand mult.} &
\textbf{Max flows} & \textbf{Target $\rho_{\mathrm{hot}}$} \\
\midrule
Low    & 8   & 400   & 80  & $<0.25$ \\
Medium & 35  & 1{,}500 & 220 & 0.40--0.60 \\
High   & 80  & 900   & 320 & 0.70--0.90 \\
Stress & 160 & 2{,}800 & 500 & $\to 1.0$ \\
\bottomrule
\end{tabular}
\end{table}

%% ============================================================
\chapter{Experiment Setup}
\label{ch:expsetup}
%% ============================================================

\begin{itemize}
  \item \textbf{Constellation:} full Walker-delta (4{,}080 satellites + 250 gateways).
  \item \textbf{Epochs:} 10 frozen-time snapshots, 9-minute steps (81 minutes total).
  \item \textbf{Pairs:} 300 random gateway-to-gateway source-destination pairs per epoch.
  \item \boldmath$d_{\mathrm{th}}$\textbf{ sweep:} 3{,}000 / 3{,}500 / 4{,}000 / 5{,}000~km.
  \item \textbf{Fixed parameters:} $\thetar = 45\text{\textdegree}$, $\thetas = 15\text{\textdegree}$.
  \item \textbf{Queuing:} 1~Gbps ISL, 200~Mbps offered load ($\rho = 0.2$),
        $L = 1{,}500$~bytes. At this load, $\Tser \approx 0.012$~ms and
        $\Tq \approx 0.0015$~ms per hop, so propagation dominates.
\end{itemize}

%% ============================================================
\chapter{Results}
\label{ch:results}
%% ============================================================

\section{Interruption Probability}

\begin{table}[H]
\centering
\caption{Interruption probability — mean over 10 epochs $\times$ 300 pairs.}
\label{tab:mainresults}
\renewcommand{\arraystretch}{1.12}
\begin{tabular}{ccccr}
\toprule
\boldmath$d_\mathrm{th}$\,(km) &
\textbf{Greedy} & \textbf{BPP pred.} &
\textbf{Dijkstra} & \textbf{Abs.\ error} \\
\midrule
3{,}000 & 0.8743 & 0.0160 & 0.0000 & 0.8584 \\
3{,}500 & 0.8490 & 0.0144 & 0.0000 & 0.8346 \\
4{,}000 & \textbf{0.7857} & \textbf{0.0144} & \textbf{0.0000} & \textbf{0.7712} \\
5{,}000 & 0.7443 & 0.0144 & 0.0000 & 0.7299 \\
\bottomrule
\end{tabular}
\end{table}

At $\dth = 4{,}000$~km, BPP predicts only $1.44\%$ interruption while the empirical
greedy simulator measures $78.57\%$---an absolute prediction error of
\textbf{77.12 percentage points}. Dijkstra achieves \textbf{zero interruptions}
across all configurations.

\section{Hop Count (\texorpdfstring{$d_\mathrm{th} = 4{,}000$}{dth=4000}~km)}

\begin{table}[H]
\centering
\caption{Mean hop count at $\dth=4{,}000$~km (delivered routes only).}
\label{tab:hops}
\begin{tabular}{lc}
\toprule
\textbf{Algorithm} & \textbf{Mean hops} \\
\midrule
Greedy   & 8.11  \\
Dijkstra & 12.65 \\
\bottomrule
\end{tabular}
\end{table}

Dijkstra uses more hops than greedy (12.65 vs 8.11 mean) because it minimises total
delay rather than hop count---shorter hops yield lower per-hop propagation. Greedy
uses fewer hops when it succeeds, but succeeds only 21.4\% of the time.

\section{End-to-End Latency (\texorpdfstring{$d_\mathrm{th} = 4{,}000$}{dth=4000}~km)}

\begin{table}[H]
\centering
\caption{Mean E2E latency (ms) at $\dth=4{,}000$~km ($\rho=0.2$, 1{,}500-byte packets).}
\label{tab:latency}
\begin{tabular}{lc}
\toprule
\textbf{Algorithm} & \textbf{Mean latency (ms)} \\
\midrule
Greedy   & 69.87 \\
Dijkstra & 66.89 \\
\bottomrule
\end{tabular}
\end{table}

Dijkstra achieves lower mean latency (66.89~ms vs 69.87~ms). At $\rho = 0.2$
the queueing contribution is negligible ($<0.1$~ms total); propagation dominates.

%% ============================================================
\chapter{Analysis: Why BPP Fails on Real Starlink Geometry}
\label{ch:analysis}
%% ============================================================

\section{The Core Contradiction}

\begin{center}
\begin{tabular}{lc}
\toprule
\textbf{Method} & \textbf{Interruption @ } $\dth=4{,}000$~\textbf{km} \\
\midrule
BPP analytical   & 0.0144 \\
Greedy empirical & 0.7857 \\
Dijkstra         & 0.0000 \\
\bottomrule
\end{tabular}
\end{center}

The BPP formula is correct under its own assumptions: with 4{,}080 randomly placed
satellites, a valid relay almost always exists somewhere in the required directional
cone. The problem is that \emph{real Starlink satellites are not random}.

\section{Walker-Delta Creates Systematic Voids}

Walker-delta orbital planes produce angular lanes with regular gaps between them. If
the routing direction from the current node to the destination points across an
inter-lane gap, no satellite satisfies $c_1$ within $\dth$---even though the
destination is reachable globally via a different path. The BPP formula averages over
isotropic placements and cannot predict these systematic directional voids.

The result is stark: BPP sees 4{,}080 satellites at high density and predicts trivially
easy routing. The greedy algorithm, working in actual Walker-delta geometry, fails
78.57\% of the time.

\section{Dijkstra as a Diagnostic}

The fact that Dijkstra achieves zero interruptions is crucial:

\begin{itemize}
  \item The Walker-delta graph is globally well-connected ($>99.6\%$ in a single component).
  \item Every tested pair has a valid route---just not one reachable by local greedy forwarding.
  \item Routing failures are entirely due to the algorithm, not to coverage gaps.
\end{itemize}

This confirms the engineering implication: real LEO networks require link-state routing
with global topology knowledge, not geographic-greedy forwarding.

\section{BPP Error Across \texorpdfstring{$d_\mathrm{th}$}{d\_th}}

\begin{table}[H]
\centering
\caption{BPP absolute prediction error — full Walker-delta constellation.}
\label{tab:bpperr}
\begin{tabular}{cccc}
\toprule
$d_\mathrm{th}$\,(km) & \textbf{BPP pred.} & \textbf{Greedy emp.} & \textbf{Abs.\ error (pp)} \\
\midrule
3{,}000 & 0.0160 & 0.8743 & 85.84 \\
3{,}500 & 0.0144 & 0.8490 & 83.46 \\
4{,}000 & 0.0144 & 0.7857 & 77.12 \\
5{,}000 & 0.0144 & 0.7443 & 72.99 \\
\bottomrule
\end{tabular}
\end{table}

Across all $\dth$ values, BPP underestimates greedy interruption by 73--86 percentage
points ($>98\%$ relative error in all cases).

%% ============================================================
\chapter{Discussion, Limitations, and Future Work}
\label{ch:discussion}
%% ============================================================

\section{Discussion}

\begin{enumerate}
  \item \textbf{Validate before applying.} A model that works under random-placement
        assumptions may be arbitrarily wrong for structured constellations. BPP should
        not be used for quantitative routing-reliability predictions on Walker-delta
        topologies without empirical validation.

  \item \textbf{Algorithm and topology interact.} The BPP algorithm was designed for
        BPP topology. Applied to Walker-delta geometry it combines algorithmic
        limitations (local-only, no backtracking) with structural voids, producing
        78.57\% interruption where BPP predicts 1.44\%.

  \item \textbf{Topology-aware routing is essential.} Dijkstra's zero interruption
        rate shows that the same satellite mesh supports reliable routing when global
        topology knowledge is available. This strongly motivates link-state protocols
        for operational LEO networks.
\end{enumerate}

\section{Limitations}

\begin{itemize}[noitemsep]
  \item Real Starlink routing, traffic, and handover policies are proprietary.
  \item Snapshots are frozen in time; real topology changes continuously.
  \item ISL assignment is a geometric heuristic, not SpaceX's actual policy.
  \item Gateway scheduling and RF interference are not modelled.
  \item The visualiser uses a simplified, calibrated traffic model---not a Starlink SLA predictor.
  \item Dijkstra is a baseline; it is not a claim about SpaceX's actual protocol.
\end{itemize}

\section{Future Work}

\begin{itemize}[noitemsep]
  \item Realistic traffic matrices (gravity models from population density).
  \item Additional routing algorithms (GPSR void avoidance, segment routing).
  \item Handover and link-failure modelling.
  \item Full parameter sensitivity sweep ($\thetar$, $\thetas$, all Starlink shells).
  \item NS-3 validation (a stub C++ scenario is already included).
  \item High-load batch experiments to quantify joint queueing + interruption effects.
\end{itemize}

%% ============================================================
\chapter{Conclusion}
\label{ch:conclusion}
%% ============================================================

This project built a complete Starlink-like multi-hop LEO routing simulator from real
TLE data. The simulator generates multi-epoch network snapshots, implements both the
Wang/BPP greedy algorithm and Dijkstra routing, and computes the closed-form BPP
analytical prediction on the same topology.

The central finding is stark. On the full 4{,}080-satellite Walker-delta constellation
at $\dth = 4{,}000$~km:

\begin{center}
\begin{tabular}{lc}
\toprule
BPP analytical prediction   & 1.44\% interruption \\
Empirical greedy routing    & 78.57\% interruption \\
Absolute prediction error   & \textbf{77.12 percentage points} \\
Dijkstra interruption       & \textbf{0.00\%} \\
\bottomrule
\end{tabular}
\end{center}

The BPP model is mathematically valid under its random-placement assumption. It fails
on real Starlink geometry because Walker-delta orbital lanes create systematic
directional voids that make greedy routing fail even where global routing succeeds.
The zero Dijkstra interruption rate confirms that the issue is the routing algorithm,
not the topology.

The practical lesson is clear: real LEO networks need link-state routing with full
topology knowledge. Geography-based greedy forwarding, as formalised in the BPP
model, is unreliable on structured Walker-delta constellations regardless of satellite
density.

The 3D visualiser makes these dynamics directly observable, showing congestion,
routing voids, and the contrast between Dijkstra paths and failed greedy routes in
real time.

%% ── Bibliography ─────────────────────────────────────────────
\bibliographystyle{IEEEtran}
\begin{thebibliography}{9}

\bibitem{wang2023reliability}
Y. Wang, M.\,A. Kishk, and M.-S. Alouini,
``Reliability Analysis of Multi-hop Routing in Multi-tier LEO Satellite Networks,''
\textit{arXiv preprint arXiv:2303.02286}, 2023.

\bibitem{bhattacherjee2019}
D. Bhattacherjee and A. Singla,
``Network topology design at 27,000 km/hour,''
in \textit{Proc. ACM CoNEXT}, 2019.

\bibitem{del2019}
E. del Portillo, B.\,G. Cameron, and E.\,F. Crawley,
``A Technical Comparison of Three LEO Satellite Network Constellations,''
in \textit{Proc. IAC}, 2019.

\bibitem{walker1984}
J.\,G. Walker, ``Satellite Constellations,''
\textit{J. British Interplanetary Soc.}, vol.~37, pp.~559--572, 1984.

\bibitem{dijkstra1959}
E.\,W. Dijkstra,
``A Note on Two Problems in Connexion with Graphs,''
\textit{Numerische Mathematik}, vol.~1, pp.~269--271, 1959.

\bibitem{skyfield}
B. Rhodes, ``Skyfield: Elegant Astronomy for Python.''
\url{https://rhodesmill.org/skyfield/}

\bibitem{threejs}
R. Cabello et al., ``Three.js -- JavaScript 3D Library.''
\url{https://threejs.org}

\bibitem{kleinrock1975}
L. Kleinrock, \textit{Queueing Systems, Vol.~I}.
Wiley, 1975.

\end{thebibliography}

\end{document}
